% Generated draft: requires mathematical and source review.
\documentclass[11pt]{article}
\usepackage[margin=25mm]{geometry}
\usepackage{fontspec}
\setmainfont{Latin Modern Roman}
\usepackage{amsmath,amssymb,mathtools,unicode-math}
\setmathfont{Latin Modern Math}
\usepackage{enumitem,needspace,xurl,hyperref}
\setlength{\parindent}{0pt}
\setlength{\parskip}{5pt}
\setlength{\emergencystretch}{3em}
\newcommand{\N}{\mathbb{N}}
\newcommand{\Z}{\mathbb{Z}}
\newcommand{\Q}{\mathbb{Q}}
\newcommand{\R}{\mathbb{R}}
\newcommand{\C}{\mathbb{C}}
\newcommand{\F}{\mathbb{F}}
\newcommand{\A}{\mathbb{A}}
\newcommand{\PP}{\mathbb{P}}
\newcommand{\E}{\mathbb{E}}
\newcommand{\eps}{\varepsilon}
\newcommand{\dd}{\,\mathrm{d}}
\begin{document}
\section*{UnsolvedMath ID: 30002211}
\textbf{Source problem code:} catalogue entry
\subsection{Definitions and mathematical statement}
Let \(D\subseteq\R\) denote the unspecified one-dimensional domain of the target density \(f:D\to[0,\infty)\), with \(\int_D f(x)\,\dd x=1\). Write \(\widehat f_\lambda\) for a regularized reconstruction, where \(\lambda\) denotes its regularization choice. The record does not specify the observation model, convolution kernel, sampling assumptions, admissible densities, or reconstruction procedure.

For clarity, one possible convention calls \(m\in\operatorname{int}(D)\) a strict local mode of a function \(u\) if
\[
\exists\delta>0\quad\forall x\in D,\qquad 0<|x-m|<\delta\ \Longrightarrow\ u(x)<u(m).
\]
This convention is illustrative: the source does not specify strictness, treatment of plateaus, boundary modes, or the regularity needed for pointwise modes of a density to be well defined independently of its almost-everywhere representative.

\textbf{Source question.} In one-dimensional nonparametric density deconvolution, suppose a regularized reconstruction exhibits an apparent additional mode that may have been introduced by regularization. Can statistical inference determine whether this feature corresponds to a genuine mode of the underlying density and whether its significance warrants rejecting the smooth reconstruction?

The source does not define the comparison smooth reconstruction, how a reconstructed mode is matched to a true mode, what makes a mode additional, or what rejection means statistically. In particular, it supplies no null and alternative hypotheses, significance level, effect-size criterion, or error guarantee. Consequently, this is an unresolved formulation of a statistical research question, rather than a fully specified testing problem. No universal quantifier over observation models or regularization methods is justified by the supplied statement.
\subsection{Short English statement}
Can one determine whether an extra peak in a regularized one-dimensional density deconvolution is a genuine feature of the underlying density, with sufficient statistical evidence to reject a smoother reconstruction? The supplied record leaves the model, mode convention, comparison, and significance criterion unspecified.
\subsection{Sources}
\begin{itemize}
\item Local UnsolvedMath registry, supplied complete source record: permanent ID \(30002211\), title \emph{Detecting Genuine Modes in Density Deconvolution}, source problem code \texttt{catalogue entry}, collection \emph{Oberwolfach Reports Open Problems}. The proposed year is recorded as 2012; no proposer is named.
\item Supplied source citation: \emph{Arbeitsgemeinschaft: Ergodic Theory and Combinatorial Number Theory} (2013), \emph{Oberwolfach Reports}, pp.\ 3061--3127, DOI 10.4171/owr/2012/51. \url{https://doi.org/10.4171/owr/2012/51}. The record describes these reports as documenting mathematical workshops and their open-problem discussions.
\item Supplied literature source, labeled ``multiscale shape inference in deconvolution'': \url{https://arxiv.org/abs/1107.1404}.
\item Supplied literature source, labeled ``original OWR report'': \url{https://doi.org/10.4171/owr/2012/51}.
\item Literature assessment dated exactly \texttt{2026-08-21}: the supplied curation labels the problem \texttt{partially\_solved}. It reports that multiscale confidence procedures can assess modes and qualitative features in one-dimensional deconvolution, while a general interpretability and significance theory for arbitrary regularized inverse reconstructions is unavailable. This is supplied dated curation derived from the OWR curation, not an independently verified assessment; its time-sensitive status and cited claims require independent verification.
\item The registry labels the statement \texttt{exact} and reports that its mathematical content was checked against the official Oberwolfach report, with notation, context, and question form normalized without changing the source-level claim. This verification assertion is supplied provenance and has not been independently checked here.
\end{itemize}
\end{document}
